% ============================================================
%  Section 4 (Results) prose + Section 5 (Conclusion), short version.
%  Labels assumed: fig:comparison, fig:beta_sweep, fig:when,
%                  fig:training_grid, fig:characteristics.
% ============================================================


% ---- end of 4.1 "Setup and baselines" ----------------------

\paragraph{The termination box.}
An episode ends when an admissible stencil for $z^{*}$ can be drawn from the
$R$-cell box $\mathcal{B}_R(z^{*})$ of \cref{sec:mdp_mask}. $R$ is a property of
the run: the box is drawn in \cref{fig:comparison} ($R=2$) and \cref{fig:when}
($R=4$), and absent from \cref{fig:beta_sweep}, \cref{fig:training_grid} and
\cref{fig:characteristics}, whose runs require a walker to arrive within half a
cell of $z^{*}$ --- the $R \to 0$ limit of the same rule. At these horizons the
two tests can be read as the same picture. All episodes are short ($t^{*}$ is
$5$, $10$ or $20$ time levels), and a walker advances one level and at most one
cell per step, so the cells from which $z^{*}$ is still reachable are few by the
time the walkers are near it: in every converged panel the terminal stencil sits
one or two cells from the star, well inside a box of half-width $2$ or $4$. The
box is an escape valve for walkers that cannot land exactly on the query cell,
not a licence to stop early and far away. Where it is drawn, picture it shrunk to
the walkers' own footprint.


% ---- append to 4.2 -----------------------------------------

\Cref{fig:training_grid} answers the question qualitatively, one deterministic
rollout per panel at three stages of training. Early policies scatter
evaluations across the domain, place points that cannot influence $z^{*}$, and
either time out or arrive expensively: $52$ evaluations for diffusion and $100$
for advection at $0.2$M steps. They then narrow --- diffusion to $34$ at $1.0$M
and $22$ at $3.1$M, advection to $24$ at $2.1$M and $15$ at $3.0$M, Burgers from
$317$ at $0.1$M to $136$ at $13.9$M and $71$ at $17.6$M. The shape says more than
the count: the converged panels commit a few walkers to a narrow column beneath
the query and leave the rest of the initial condition untouched. Burgers also
shows the price --- an order of magnitude more environment steps than the other
two, and a converged rollout four times more expensive at twice the horizon.

\Cref{fig:comparison} makes the comparison quantitative for a single query on
diffusion at $t^{*}=10\,\Delta t$. The uniform march must fill the FTCS domain of
dependence of $z^{*}$, $100$ evaluations for $6.4\times10^{-7}$ (nodes outside
the pyramid are drawn but not charged, so this is the strongest honest form of
the classical baseline); masked-random spends $350$ for $8.3\times10^{-6}$; the
policy spends $22$ for $3.4\times10^{-5}$. The error column must be read with the
reward in mind: this run was trained at $e_{\mathrm{tol}}=10^{-3}$, so all three
errors give $e_{\mathrm{norm}}=0$ and the accuracy term in \eqref{eq:reward}
contributes nothing. Cost is the only live axis, and there the policy uses
$4.5\times$ fewer evaluations than the uniform march and $16\times$ fewer than
masked-random, three decades inside the tolerance it was asked to respect. It is
a Pareto point, not a domination: the policy is cheaper and less accurate because
that is what it was paid to be.


% ------------------------------------------------------------
\subsection{Accuracy is bought by geometry, not by more points}
\label{sec:results_beta}
% ------------------------------------------------------------

\Cref{fig:beta_sweep} makes the accuracy term active and holds everything else
fixed, sweeping $\beta \in \{10^{3},10^{4},10^{5}\}$ in the original linear
reward $R=-(\beta|e|+\lambda k)$ at $\lambda=0.01$. The reading is marginal: an
extra evaluation costs $\lambda$ and pays for itself only while it buys more than
$\lambda/\beta$ of error, so the policy drives the error to
$\varepsilon_{\mathrm{eff}}=\lambda/\beta$ and stops. The sweep is therefore over
an effective tolerance, $\varepsilon_{\mathrm{eff}}=10^{-5},10^{-6},10^{-7}$,
playing the role $e_{\mathrm{tol}}$ plays in the current reward.

The achieved error tracks it --- $1.0\times10^{-4}$, $3.8\times10^{-6}$,
$5.3\times10^{-7}$, more than two decades --- while the cost does not:
$k=13,15,11$, with no trend. A classical scheme buys accuracy by refining, and
the naive expectation is that an agent paid for accuracy would place more points.
It does not; at fixed cost it rearranges the points it already has. With
$\Delta x \in \{-1,0,+1\}$ cells and $\Delta t \in \{1,2\}$ levels, some walkers
step one level and some two, so the terminal stencil reaches $z^{*}$ from several
time levels at once. That is a geometry, not a budget, and it is the sense in
which the object being adapted is the scheme rather than the resolution: the
policy can change the stencil shape at constant point count, and
\eqref{eq:bound_pde} says why that pays, since $\|\mathbf{w}\|_1$ and $h$ depend
on how the neighbours are arranged and not on how many there are.


% ------------------------------------------------------------
\subsection{Placement is also \emph{when}}
\label{sec:results_when}
% ------------------------------------------------------------

The action space contains a third decision that is easy to overlook: the order in
which points are computed. Every interior value is a reconstruction from
neighbours that already exist, and moves are solved non-causally, so a point may
be built from neighbours both below and above it in time. \Cref{fig:when} shows
the same eleven points, the same three walker chains and the same five-point
stencil at $z^{*}$ in three construction orders. Two are the policy's own, at the
$6.0$M and $7.2$M checkpoints of one advection run, differing by two adjacent
transpositions; their errors are $8.9\times10^{-5}$ and $7.7\times10^{-5}$, and
the third order gives $3.5\times10^{-5}$. The point set fixes the stencil and
hence the weights, so the error is $|\sum_i w_i u_i - u(z^{*})|$ at fixed
$\mathbf{w}$: the entire spread comes from the computed $u_i$, that is, from the
order alone.

Enumerating all $11550$ interleavings that preserve each chain's internal order
gives errors from $9.5\times10^{-8}$ to $1.2\times10^{-3}$, median
$1.2\times10^{-4}$ --- a factor of $1.3\times10^{4}$ on a fixed cloud. The
policy's own orders sit near the median, not the minimum, which is the honest
reading: the reward is terminal, the order is decided one step at a time without
observing the field, and nothing in the formulation tells the agent this axis
exists. Four decades are available there at no cost in evaluations.


% ------------------------------------------------------------
\subsection{What geometry does the policy discover?}
\label{sec:results_geometry}
% ------------------------------------------------------------

If the rule reflected the equation rather than the shortest route to the query,
the cloud should lean along the direction information travels: slope $c=0.5$ for
advection, the solution itself for Burgers. \Cref{fig:characteristics} sets the
converged rollouts against the characteristic field, and the half-width of the
visited cloud against lookback $\tau=(t^{*}-t)/\Delta t$ against the three scales
that could set it --- characteristic $c\,\tau\,\Delta t/\Delta x$, diffusive
$\sqrt{\alpha\tau\Delta t}/\Delta x$, and the action cone $\tau$, one cell per
level.

The answer is the action cone: the Burgers envelope is $w=\tau+1$ cells for every
$2\le\tau\le19$ and the advection envelope is $w=\tau$ exactly for
$2\le\tau\le4$, the widest structure the action set permits. We report that as a
negative result, but it is not evidence that the policy ignores the equation. Both
runs take $\Delta t$ from the diffusive stability limit, giving advective Courant
numbers of $0.23$ and $0.046$, so over the whole lookback the characteristic
through $z^{*}$ drifts $1.14$ and $0.92$ cells --- less than the one cell a
walker moves in a single step. Following it is not an available behaviour, and
the figure rules the explanation out rather than testing it. The one scale in
range is the diffusive one, about $3$ cells at $t^{*}$ for Burgers and $1.5$ for
advection; the stencil at $z^{*}$ does live on it even though the cloud feeding
it does not, consistent with \eqref{eq:bound_pde}, where the bound is set by the
stencil radius at the query and not by the shape of the supporting cloud. Testing
whether a placement rule can recover a transport direction needs a run at Courant
number of order one, which this record does not contain.


% ---- replacement caption for fig_characteristics -----------
% \caption{Does the placement rule follow the characteristic? Left: converged
% rollouts with the characteristic field through them --- slope $c$ for
% advection, and for Burgers the curves from integrating $dx/dt=u(x,t)$ on the
% Hopf--Cole reference. Right: half-width of the visited cloud against lookback
% $\tau=(t^{*}-t)/\Delta t$ against the three scales that could set it. The cloud
% sits on the action cone, one cell per level, in both rows; the characteristic
% drifts only $1.14$ (advection) and $0.92$ (Burgers) cells over the whole
% lookback, because $\Delta t$ is fixed by the diffusive limit. The action cone
% is the set of points a walker could still reach $z^{*}$ from, not a bound on
% where a walker may sit, which is why the $\tau=1$ marker lies above it.}


% ============================================================
\section{Conclusion}
\label{sec:conclusion}
% ============================================================

We have treated discretisation as a sequential decision problem rather than a
choice fixed before the solve begins. The two halves meet at one interface. The
reconstruction half is classical: a Taylor-constrained GFDM solve returns
minimum-norm weights for any scattered space--time neighbourhood, with the PDE
substituted into the expansion so the equation supplies constraints rather than
merely being approximated --- which raises the local order from $h^{2}$ to
$h^{3}$ under parabolic scaling and reproduces the forward-Euler and fourth-order
stencils to machine precision on a regular grid. It leaves no freedom: once the
points are chosen, the weights follow. The learning half therefore has a clean
job, choosing the points, and it is the only half that is learned; the agent
never sees the field, only geometry relative to the query, and is scored on a
terminal trade-off between error and evaluations.

Trained with maskable PPO on diffusion, advection--diffusion and viscous Burgers,
the policy learns a placement rule: it narrows from scattering evaluations to
committing a few walkers to a column beneath the query, answering one diffusion
query three decades inside tolerance in $22$ evaluations against $100$ for the
minimum-cost uniform march and $350$ for masked-random. Tightening the tolerance
by two decades drives the achieved error down while the evaluation count stays
flat: accuracy is bought by rearranging points, not by adding them, through
stencils that reach the query from several time levels at once. That is what
distinguishes adapting the scheme from adapting the resolution.

The record's limits are equally clear. The converged geometry is the action cone,
not a physical scale --- though at these Courant numbers the characteristic
drifts less than one cell over the whole horizon, so no transport-following
behaviour was available to be learned. The construction order, which the action
space already controls, spans four decades of terminal error on a fixed cloud,
and the learned orders sit near the median of that spread: a lever the
formulation exposes but the reward does not reach. Training cost is not
competitive with solving cost, Burgers needing on the order of $10^{7}$
environment steps for one query, and every result here is a single query in one
dimension. Each is a direction as much as a caveat. Amortising a policy over many
queries would make the training cost accountable, and the state --- displacements
in grid cells relative to the query --- was formulated to be reusable across
queries and resolutions for that reason. A run at Courant number of order one
would pose the geometry question this lattice cannot. Rewarding the construction
order closes the largest unexploited gap we found. And higher dimensions are
where the argument should pay: uniform stencils grow with dimension, the GFDM
construction does not care about it, and which few points to place is exactly the
question that becomes harder to answer by hand.
